\documentclass[10pt,a4paper]{amsart}
\usepackage[margin=19mm]{geometry}
\usepackage{amsmath,amssymb,mathtools}
\usepackage[scr=rsfs,cal=euler]{mathalfa}
\usepackage{xfrac}
\usepackage[unicode,hidelinks,bookmarks=false]{hyperref}
\newcommand{\ie}{i.e.\ }
\newcommand{\cf}{cf.\ }
\newcommand{\resp}{respectively\ }
\newcommand{\Subsection}{\subsection}
\providecommand{\nobreakdash}{\nobreak\hbox{-}}
\setlength{\emergencystretch}{2em}
\multlinegap0pt
\usepackage[shortlabels]{enumitem}
\usepackage{amssymb}
\usepackage{float}
\newtheorem{lemma}{Lemma}
\newcommand{\E}{\mathbb{E}}
\newcommand{\N}{\mathbb{N}}
\renewcommand{\P}{\mathbb{P}}
\newcommand{\Z}{\mathbb{Z}}
\newcommand{\cF}{\mathcal F}
\newcommand{\cG}{\mathcal G}
\newcommand{\del}[1]{{\color{blue} #1}}
\title{Total disconnectedness and percolation for the supports of super-tree random measures}
\author{Edwin Perkins, Delphin S\'enizergues}
\date{Source: arXiv:2409.11841v2 (CC BY 4.0)}
\begin{document}
\maketitle

\noindent\emph{Source and licence.} Total disconnectedness and percolation for the supports of super-tree random measures. Version arXiv:2409.11841v2 (CC BY 4.0), distributed by arXiv under the Creative Commons Attribution 4.0 International licence, http://creativecommons.org/licenses/by/4.0/, as recorded in the arXiv licence metadata for this exact version and shown on the abstract page https://arxiv.org/abs/2409.11841v2. Full licence notice: \texttt{licenses/arxiv-2409.11841-CC-BY-4.0.txt}.\par\smallskip

\noindent\textit{Editorial scope.} Counted unit: the article's lemma gammaext (source0.tex lines 1618--1626). The article's complete original proof of it is delivered with it (source0.tex lines 1699--1827, one \texttt{proof} environment of 129 lines, which the article places away from the statement). No mathematics is written, repaired or re-derived: the statement and the proof are the article's own lines, in the article's own notation, and no retained line is substituted or re-flowed.  Retained uncounted as the source-local context the proof needs: lines 1335--1337; lines 1525--1530; lines 1558--1565; lines 1567--1571; lines 1686--1690; lines 1687--1688.  Nothing was omitted from any retained span, so the package carries no omission record.  No retained line contains a citation, so the wrapper carries no bibliography and no bibliography entry.  No retained line contains a figure environment, an image-inclusion command or drawing code, so no image is delivered and the package declares no asset dependency.  Editorial formatting only, in addition: an AMS \texttt{amsart} preamble that restates the article's own declarations and macros used by the retained lines, namely the environments \texttt{lemma} on separate counters, together with the standard packages those lines need.  The retained environments print in this document as Lemma 1 in order of appearance, whereas the article prints the same items with the numbering of its own full document.  The complete native source of the article as distributed in the arXiv e-print for this exact version is bundled unchanged as \texttt{sources/165443/source0.tex}.
\begin{equation}\label{veenotation}
\text{$g\vee \ell\in F_{m+1}$ be the index obtained by adding $\ell$ to $g$ as the last digit.}
\end{equation} 

\begin{lemma}\label{ellneighb} Let $x=(x^1,\dots,x^d), y=(y^1,\dots,y^d)\in G_m$, for $m\in\N$. If $L=L(x,y)$, then
\[\text{$C_m(x)$ and $C_m(y)$ are neighbours $\iff$ $x\neq y$ and for any $k\in L^c$, $|x^k-y^k|=B^{-m}$.}\]
In this case
\begin{equation}\label{Cinters}\overline{C_m(x)}\cap\overline{C_m(y)}=\{z:z^k\in[x^k,x^k+B^{-m}]\text{ for }k\in L,\text{ and }z^k=x^k\vee y^k\text{ for }k\in L^c\}.
\end{equation}
\end{lemma}

\begin{lemma}\label{ellnab}
Let $n_1\le n_2$  be in $\N$, $f,g\in F_{n_2}$, and $\ell\in\{0,\dots,d-1\}$.
%\del{The following properties hold deterministically.}
\begin{enumerate}[label=\emph{(\alph*)}, ref=(\alph*)]
\item\label{it:L is non-increasing} $L(S^f,S^g)\subset L(S^{f|n_1},S^{g|n_1})$, and so in particular, $S^{f|n_1}\neq S^{g|n_1}$ implies $S^f\neq S^g$.
\item\label{it:neighbours descend from neighbours} If $C_{n_2}(S^f)$ and $C_{n_2}(S^g)$ are $\ell$-neighbours and $S^{f|n_1}\neq S^{g|n_1}$, then $C_{n_1}(S^{f|n_1})$ and $C_{n_1}(S^{g|n_1})$ are $\ell'$-neighbours for some $\ell'\ge \ell$.
\end{enumerate}
\end{lemma}

\begin{align}\label{Sincr}
\nonumber|S^{f,k}-S^{g,k}|&\ge |S^{f|n_1,k}-S^{g|n_1,k}|-\sum_{m=n_1+1}^{n_2}B^{-m}|X^{f|m,k}-X^{g|m,k}|\\
\nonumber&\ge |S^{f|n_1,k}-S^{g|n_1,k}|-B^{-n_1-1}(B-1)\sum_{m=0}^{n_2-n_1-1}B^{-m}\\
&= |S^{f|n_1,k}-S^{g|n_1,k}|-B^{-n_1}+B^{-n_2}.
\end{align}

\begin{lemma}\label{supermart} Assume $\{M_n:n\in\N\}$ is a $\Z_+$-valued $(\cF_n)$-supermartingale such that for any $k\in\N$ there is a $r_k>0$ so that for all $n\in\N$,
\begin{equation}\label{martcond}\P(M_{n+1}\neq k \, | \, \cF_n)\ge r_k\text{ on }\{M_n=k\}.
\end{equation}
Then $M_n\to 0$ a.s.
\end{lemma}

\begin{equation}\P(M_{n+1}\neq k \, | \, \cF_n)\ge r_k\text{ on }\{M_n=k\}.
\end{equation}

\begin{lemma}\label{gammaext}
Assume $d\ge\frac{4}{\beta}-1$, $\ell\in\{0,\dots,d-1\}$, and $R\in\N\cup\{\infty\}$ is a random
variable.
% vector.} 
With probability $1$, 
if $R<\infty$ and $f,g\in F_R$ are such that $C_R(S^f)$ and $C_R(S^g)$ are not $\ell'$-neighbours for any $\ell'\in\{\ell+1,\dots,d-1\}$, then 
\begin{equation}\label{Rfinite}R_R^\ell(f,g):=\inf\{n\ge R:\, \Gamma^\ell_{R,n}(f,g)=\emptyset\}<\infty.
\end{equation}
\end{lemma}

\begin{proof}[Proof of Lemma~\ref{gammaext}]
As noted after the statement of Lemma~\ref{gammaext}, we may assume $f$ and $g$ are $\ell$-neighbours when verifying \eqref{Rfinite}. 
Next we claim that it suffices to prove the result for $R=m$ a constant.
Assume this constant case and let $f,g\in F_R$ be $\ell$-neighbours.  
So on $\{R=m\}$, $f,g\in F_m$ are $\ell$-neighbours. 
The result for $R$ identically $m$ gives w.p.$1$ $R_m^\ell(f,g)<\infty$.  
This gives \eqref{Rfinite} w.p.$1$ on $\{R=m\}$.  
Combine the resulting null sets to conclude that w.p.$1$, \eqref{Rfinite} holds.  

So we clearly may fix $m\in\N$, $f\neq g\in F_m$ and $\ell\in\{0,\dots,d-1\}$, and it suffices to prove that with probability $1$ on 
\[\Omega_\ell=\{\text{$f$, $g$ are $\ell$-neighbours}\}\]
which is in $\cG_m$, we have
\begin{equation}\label{Rmfinite}
R^\ell_m(f,g)<\infty.
\end{equation}
Until otherwise indicated we assume $\omega\in\Omega_\ell$. 
We simplify notation and for $n\ge m$ write $\Gamma_n^{\ell'}(f,g)$ for $\Gamma_{m,n}^{\ell'}(f,g)$. 

We first claim that 
\begin{equation}\label{constL}
\forall n\ge m, \forall (f',g')\in \Gamma_n^\ell(f,g), \quad L(S^{f'},S^{g'})=L(S^f,S^g).
\end{equation}
For this, note that Lemma~\ref{ellnab}\ref{it:L is non-increasing} implies that for all $n\ge m$ and $(f',g')\in K^f_n\times K_n^g\supset \Gamma_n^{\ell}(f,g)$, we have  $L(S^{f'},S^{g'})\subset L(S^f,S^g)$. 
On the other hand  $(f',g')\in\Gamma^\ell_n(f,g)$ implies $|L(S^{f'},S^{g'})|=\ell=|L(S^f,S^g)|$ (recall we work on $\Omega_\ell$), and so \eqref{constL} follows. 

If $n\ge m$ and $x,y\in G_n$, let 
\[L'_n(x,y)=\{1\le k \le d:x^k=y^k+B^{-n}\}.\]
It follows from Lemma~\ref{ellneighb} that ($\dot\cup$ denotes a disjoint union)
\begin{equation}\label{Lcdecomp}
\text{if $f',g'\in F_n$ are neighbours then $L(S^{f'},S^{g'})^c=L'_n(S^{f'},S^{g'})\dot\cup L'_n(S^{g'},S^{f'})$.}
\end{equation}
Together \eqref{constL} and \eqref{Lcdecomp} give us
\begin{equation}\label{unioneq}
L'_n(S^{f'},S^{g'})\cup L'_n(S^{g'},S^{f'})=L'_m(S^f,S^g)\cup L'_m(S^g,S^f)\ \forall (f',g')\in \Gamma_n^\ell(f,g)\ \forall \ n\ge m.
\end{equation}
Next we claim that 
\begin{equation}\label{constL'}
\forall n\ge m, \forall(f',g')\in \Gamma_n^\ell(f,g), \ L'_n(S^{f'},S^{g'})=L'_m(S^f,S^g)\text{ and }L'_n(S^{g'},S^{f'})=L'_m(S^g,S^f).
\end{equation}
For this, argue as in \eqref{Sincr} but without the absolute values, to see that for $(f',g')\in\Gamma^\ell_n(f,g)$ $(n\ge m$) and $k\in L'_m(S^f,S^g)$, 
\begin{align*}
 S^{f',k}-S^{g',k}&=S^{f,k}-S^{g,k}+\sum_{j=m+1}^n B^{-j}(X^{f'|j,k}-X^{g'|j,k})\\
&=B^{-m}+\sum_{j=m+1}^nB^{-j}(X^{f'|j,k}-X^{g'|j,k})\\
&\ge B^{-n}.
\end{align*}
The fact that $f'$ and $g'$ are neighbours implies $|S^{f',k}-S^{g',k}|=0$ or $B^{-n}$ by Lemma~\ref{ellneighb}. So the above implies $S^{f',k}-S^{g',k}=B^{-n}$, and we have proved $L'_m(S^f,S^g)\subset L'_n(S^{f'},S^{g'})$, and therefore $L'_m(S^g,S^f)\subset L'_n(S^{g'},S^{f'})$ by symmetry.  
These inclusions and the equality in \eqref{unioneq} now give \eqref{constL'}.

We claim that the evolution in $n\ge m$ of  $\Gamma^\ell_n(f,g)$ satisfies
\begin{align}\label{gammaevol}
(f',g')\in \Gamma^\ell_{n+1}(f,g)\ \ \ \text{implies }&(\pi f',\pi g')\in\Gamma^\ell_n(f,g),\ f'_{n+1}\le Z^{\pi f'},\ g'_{n+1}\le Z^{\pi g'},\\ \nonumber&\forall k\in L(S^f,S^g), \ X^{f',k}=X^{g',k},\\
\nonumber&\forall k\in L'_m(S^f,S^g), \ X^{g',k}=B-1, X^{f',k}=0,\text{ and}\\
\nonumber&\forall k\in L'_m(S^g,S^f), \ X^{f',k}=B-1, X^{g',k}=0.
\end{align}
(In fact it is also not hard to show the converse implication holds in the above but we will not need this.)  To prove the above, note first that for any $n\ge m$ we clearly have
\begin{equation}\label{Kevol}
(f',g')\in K^f_{n+1}\times K^f_{n+1}\quad \text{iff} \quad (\pi f',\pi g')\in K^f_n\times K^g_n \text{ and } f'_{n+1}\le Z^{\pi f'}\text{ and }g'_{n+1}\le Z^{\pi g'}.
\end{equation}
Let $(f',g')\in \Gamma^\ell_{n+1}(f,g)$ where $n\ge m$. Lemma~\ref{ellnab}\ref{it:L is non-increasing} implies
\[L(S^{f'},S^{g'})\subset L(S^{\pi f'},S^{\pi g'})\subset L(S^f,S^g).\]
The cardinalities of the first and last sets above are both $\ell$ (recall we are on $\Omega_\ell$), and so we conclude
\begin{equation}\label{Lsame}
L(S^{f'},S^{g'})=L(S^{\pi f'},S^{\pi g'})=L(S^f,S^g),
\end{equation}
and
\begin{equation}\label{Lcard}
|L(S^{\pi f'},S^{\pi g'})|=\ell.
\end{equation}
Recall that $S^f\neq S^g$ (they are in fact $\ell$-neighbours on $\Omega_\ell$) and so by Lemma~\ref{ellnab}\ref{it:L is non-increasing}, $S^{\pi f'}\neq S^{\pi g'}$. 
We can therefore apply Lemma~\ref{ellnab}\ref{it:neighbours descend from neighbours} to see that $\pi f'$ and $\pi g'$ are $\ell'$-neighbours for some $\ell'\ge \ell$, and hence are neighbours, because $f'$ and $g'$ are $\ell$-neighbours.  
Now \eqref{Lcard} shows that $\pi f'$ and $\pi g'$ must be $\ell$-neighbours.  
This conclusion and \eqref{Kevol} show that 
\begin{equation}\label{piinG}
(\pi f',\pi g')\in \Gamma_n^\ell(f,g).
\end{equation}
By \eqref{Lsame} for all $k\in L(S^f,S^g)$, $S^{\pi f',k}=S^{\pi g',k}$ and $S^{f',k}=S^{g',k}$, which implies
\begin{equation}\label{equalonL}
\forall k\in L(S^f,S^g), \ X^{f',k}=X^{g',k}.
\end{equation}
By \eqref{piinG} we may apply \eqref{constL'} to $(\pi f',\pi g')$, as well as $(f',g')$, and deduce
\begin{equation*}
L'_n(S^{\pi f'},S^{\pi g'})=L'_m(S^f,S^g)\quad \text{and} \quad L'_{n+1}(S^{f'},S^{g'})=L'_m(S^f,S^g).
\end{equation*}
By definition this means
\[\forall k\in L_m'(S^f,S^g), \quad S^{\pi f',k}=S^{\pi g',k}+B^{-n}\text{ and }S^{f',k}=S^{g',k}+B^{-n-1}.
\]
Take differences in the above to conclude that $X^{f',k}=X^{g',k}+1-B$, and therefore we obtain
\begin{equation}\label{L'cond2}
\forall k\in L'_m(S^f,S^g), \ X^{g',k}=B-1\text{ and }X^{f',k}=0.
\end{equation}
Reversing the roles of $f$ and $g$ gives
\begin{equation}\label{L'cond3}
\forall k\in L'_m(S^g,S^f), \ X^{f',k}=B-1\text{ and }X^{g',k}=0.
\end{equation}
Now combine \eqref{Kevol}, \eqref{piinG}, \eqref{equalonL}, \eqref{L'cond2} and \eqref{L'cond3} to complete the proof of \eqref{gammaevol}. 

Recall from \eqref{veenotation} the notation $f'\vee k\in F_{n+1}$ if $f'\in F_n$ and $k\in\N$. For $n\ge m$ define $M_{n}=1_{\Omega_\ell}|\Gamma_n^\ell(f,g)|$. Clearly $M_n$ is $\cG_n$-measurable. By \eqref{gammaevol} for $n\ge m$,
\begin{align}\label{Mnbnd}
M_{n+1}\le 1_{\Omega_\ell}\sum_{(f',g')\in \Gamma_n^{\ell}(f,g)}\sum_{i=1}^{Z^{f'}}\sum_{j=1}^{Z^{g'}}&1(\forall k\in L(S^f,S^g), \ X^{f'\vee i,k}=X^{g'\vee j,k})\\
\nonumber&\times 1(\forall k\in L_m'(S^f,S^g), \ X^{g'\vee j,k}=B-1,\,X^{f'\vee i,k}=0)\\
\nonumber&\times 1(\forall k\in L_m'(S^g,S^f), \ X^{f'\vee i,k}=B-1,\,X^{g'\vee j,k}=0).
\end{align}
Let $\bar\cG_n=\cG_n\vee\sigma(Z^{f'},f'\in F_n)$. 
Condition first on $\bar \cG_n$ and use the independence properties of $(X^f,f\in F)$ and $(Z^f,f\in F)$ to see that 
\begin{align*}
\E(M_{n+1}\, | \, \cG_n)&\le 1_{\Omega_\ell}\sum_{(f',g')\in \Gamma^\ell_n(f,g)}\E\Bigl(\sum_{i=1}^{Z^{f'}}\sum_{j=1}^{Z^{g'}}B^{-\ell}B^{-2(|L_m'(S^f,S^g)|+|L_m'(S^g,S^f)|)} \ \Bigl| \ \cG_n\Bigr)\\
&=M_n\mu^2B^{-\ell}B^{-2(d-\ell)}\quad\text{(by \eqref{Lcdecomp} and $|L(S^f,S^g)^c|=d-\ell$)}\\
&\le M_nB^{4/\beta}B^{-1-d}\le M_n,
\end{align*}
where we have used $\ell\le d-1$, $\mu=B^{2/\beta}$ and $d\ge 4/\beta-1$ in the last line. Therefore $(M_n)$ is a $\Z_+$-valued $(\cG_n)$-supermartingale.  

To prove that \eqref{Rmfinite} holds a.s.\ on $\Omega_\ell$ it clearly suffices to show $M_n\to 0$ a.s.\ and for this we will verify the hypothesis \eqref{martcond} of Lemma~\ref{supermart}. 
On $\Omega_\ell$ we have by \eqref{Lcdecomp} that $0<|L(S^f,S^g)^c|=|L'_m(S^f,S^g)|+|L'_m(S^g,S^f)|$. 
So one of $L_m'(S^f,S^g)$ or $L'_m(S^g,S^f)$ is nonempty and we may assume without loss of generality it is the former and set $k_m=\min L'_m(S^f,S^g)$, which is $\cG_m$-measurable.  
Let $\gamma_{n}^\ell(f,g)\subset K^f_n$ be the projection of $\Gamma_n^\ell(f,g)$ onto the first variable. 
By \eqref{Mnbnd} on $\{L'_m(S^f,S^g)\neq \emptyset\}\cap\Omega_\ell \,\in\cG_m$, 
\begin{align}\label{lbnd0}
\nonumber\P(M_{n+1}=0\, | \, \cG_n)&\ge\P\Bigl(\cap_{(f',g')\in \Gamma_n^\ell(f,g)}\cap_{i=1}^{Z^{f'}}\cap_{j=1}^{Z^{g'}}\{\exists k\in L'_m(S^f,S^g)\\
\nonumber&\phantom{\ge\P\Bigl(\cap_{(f',g')\in \Gamma_n^\ell(f,g)}\cap_{i=1}^{Z^{f'}}\cap_{j=1}^{Z^{g'}}} s.t.\ (X^{f'\vee i,k},X^{g'\vee j,k})\neq(0,B-1)\}\ \Bigl| \ \cG_n\Bigr)\\
&\ge \P\Bigl(\cap_{f'\in \gamma_{n}^\ell(f,g)}\cap_{i=1}^{Z^{f'}}\{X^{f'\vee i,k_m}\neq 0\}\ \Bigl| \ \cG_n\Bigr).
\end{align}
Conditional on $\bar\cG_n$, $\{1(X^{f'\vee i,k_m}\neq 0):f'\in\gamma^\ell_{n}(f,g),i\le Z^{f'}\}$ are i.i.d.\ $\mathrm{Bernoulli}(p)$ r.v.'s with $p=1-B^{-1}$, while conditional on $\cG_n$, $\{Z^{f'}:f'\in\gamma^\ell_{n}(f,g)\}$ are i.i.d.\ copies of the branching random variable $Z$.  
So condition first on $\bar\cG_n$ and then on $\cG_n$ to see that the right-hand side of \eqref{lbnd0} is (on $\Omega_\ell$)
\[\E\Bigl(\prod_{f'\in\gamma^\ell_{n}(f,g)}(1-B^{-1})^{Z^{f'}}\ \Bigl| \ \cG_n\Bigr)=\E\Bigl((1-B^{-1})^Z\Bigr)^{|\gamma^\ell_{n}(f,g)|}\ge \E((1-B^{-1})^Z)^{M_n},\]
where the last inequality holds because $|\gamma^\ell_{n}(f,g)|\le |\Gamma^\ell_n(f,g)|=M_n$ on $\Omega_\ell$. 
Recall that $M_{n+1}=0$ on $\Omega_\ell^c$ and so for any $k\in\N$, on $\{M_n=k\}$,
\[\P(M_{n+1}=0 \, | \, \cG_n)\ge \E((1-B^{-1})^Z)^{k}=:r_k>0.\]
This proves \eqref{martcond} and the proof is complete.
\end{proof}

\end{document}
